\documentclass{article}
\usepackage{icml2027}
\usepackage{booktabs,amsmath,amssymb,graphicx,xcolor,microtype,url}
\usepackage[hidelinks]{hyperref}

\icmltitlerunning{Transfer That Survives Its Controls}

\begin{document}

\twocolumn[
\icmltitle{Transfer That Survives Its Controls:\\
A Validity Protocol for Cross-Representation Claims in Biological Sequence Models}
\icmlsetsymbol{equal}{*}
\begin{icmlauthorlist}
\icmlauthor{Saanvi Skanda Subramanian}{caltech}
\end{icmlauthorlist}
\icmlaffiliation{caltech}{Computation and Neural Systems, California Institute of Technology}
\icmlcorrespondingauthor{Saanvi Skanda Subramanian}{ssubramanian@caltech.edu}
\icmlkeywords{evaluation, scientific machine learning, protein language models, genomic language models}
\vskip 0.3in
]

\printAffiliationsAndNotice{}

\begin{abstract}
DNA and protein are two encodings of one molecule related by a map that is exactly
known. This makes biological sequence models an unusually sharp setting for a
question the rest of machine learning can only ask loosely: when a model trained
on one representation appears to know something about another, is that transfer,
or an artifact of how the question was posed? We give four controls that decide
this, each cheap and each reusable outside biology, and apply them to genomic and
protein language models. The controls are a noise ceiling read off the model's own
violation of an invariance it should respect, a pair of manipulations that destroy
the representation map while preserving surface statistics, composition-matched
context twins, and trivial baselines scored under correlated-unit accounting.
Applying them overturns four published or prior claims, three of them our own. A
genomic representation that scored $\rho=+0.141$ on held-out residues, read as
weak transfer, scores $+0.143$ after its reading frame is destroyed and $+0.129$ on
the reverse complement. A context dose-response is reproduced in both directions by
flanks carrying no genomic information. Across 25 deep mutational scanning assays
in 9 organisms, the surviving picture is that genomic-to-protein transfer is a
property of architecture and corpus and not of scale: the Nucleotide Transformer
family is flat at $-0.013$ over four sizes, Evo~2 7B reaches $+0.266$ but is
statistically indistinguishable from BLOSUM62 at $+0.228$, and the Nucleotide
Transformer scaling trend under-predicts Evo~2 by a factor of eleven. Released checkpoints let us separate two
candidate explanations at matched architecture: codon-aligned tokenization and a
thousandfold increase in training-corpus diversity each leave the result unchanged,
with paired differences of $+0.016$ and $+0.005$ that both contain zero. We report
the cost of each control in model evaluations so that the protocol can be adopted
without a compute argument.
\end{abstract}

\section{Introduction}

Claims that a model trained on one representation has learned something about
another are common and consequential. A genomic language model that scored protein
fitness would let practitioners rank variants without a protein model, and would be
evidence that the model had internalised the genetic code and not merely nucleotide
statistics. Similar claims are made for multimodal encoders, for code models
reasoning about natural language, and for any system said to have acquired a
capability it was not trained on.

These claims are hard to check because the representations are not independent.
Under any encoding rule the codon determines the amino acid, so the nucleotides at
a substituted site already \emph{are} residue identity. A probe that recovers
residue identity from DNA has not demonstrated a protein-level representation; it
has read a lookup table. The difficulty is general. Whenever representation $A$
determines property $P$ of modality $B$ by a known map, a measurement that claims
$A$ encodes $P$ ``at the level of $B$'' must rule out the map itself.

Our contribution is a protocol that rules it out, and a case study where applying
it changes the conclusion four times. We state the protocol first because it is the
reusable part.

\paragraph{Contributions.}
\begin{enumerate}\itemsep0pt
\item \textbf{A noise ceiling obtained from the model, not the data}
(Section~\ref{sec:ceiling}). When many inputs must share a label by construction,
the spread of model scores across them bounds the correlation any proxy built on
that model could achieve. For genomic models scoring coding sequences, synonymous
encodings supply this set at no extra cost, since a single masked pass already
prices every alternative. The resulting attenuation ceiling is $0.900$, so a null
at $-0.013$ cannot be attributed to measurement noise. Without such a ceiling a
null states nothing.
\item \textbf{Map-destroying controls} (Section~\ref{sec:map}). Frameshifting and
reverse-complementing destroy the DNA-to-protein map while leaving the sequence a
plausible genomic string. A representation that is protein-level must lose its
signal under both. The claim we tested did not.
\item \textbf{Composition-matched context twins} (Section~\ref{sec:context}).
Context dose-response curves are reproduced by flanks matched on mononucleotide or
dinucleotide composition and carrying no genomic information, in both directions of
sign.
\item \textbf{Trivial baselines under correlated-unit accounting}
(Section~\ref{sec:trivial}). BLOSUM62 alone reaches 49\% of ESM-2 across the same
assays. A 60-parameter one-hot vector aligns with ESM-2 better than a genomic model
does. On the genomic models' own benchmark tasks a positional one-hot wins
outright. 194 ProteinGym assays are 165 independent proteins and must be counted
as such.
\item \textbf{The case study}: genomic-to-protein transfer is an architecture and
corpus property, not a scaling one (Section~\ref{sec:scale}).
\end{enumerate}

We report four retractions of our own earlier claims in Section~\ref{sec:retract}.
We do this in the body, not an appendix, because the protocol's value is
precisely that it caught them, and because a reader deciding whether to adopt it
should see what it costs to not have had it.

\section{Related Work}

\paragraph{Protein and genomic language models.} ESM-2 \citep{lin2023esm2} and
ESM-C established masked protein language models as zero-shot variant-effect
predictors, evaluated at scale by ProteinGym \citep{notin2023proteingym}. On the
nucleotide side, the Nucleotide Transformer \citep{dalla2024nucleotide}, HyenaDNA
\citep{nguyen2023hyenadna} and Evo~2 \citep{brixi2025evo2} model genomes directly.
Evo and Evo~2 have been used to design functional sequences
\citep{merchant2025evo,king2026phage}, though in the latter case through
inference-time guidance from separate predictive models and not bare likelihood
ranking.

\paragraph{Genomic models on protein tasks.} \citet{boshar2024glm} ask our question
directly and reach an apparently opposite answer: under supervised fine-tuning on
five protein tasks, Nucleotide Transformer v2 500M matches or outperforms ESM-2 on
three of them. They attribute the two losses to tasks needing codon-level
resolution and train a 3-mer variant to address it, which improves those two tasks
and leaves the rest unchanged. Their protocol differs from ours in the way that
matters: they fine-tune, we score zero-shot likelihood; they use no
non-model baseline, no map-destroying control and no noise ceiling; and they do not
evaluate zero-shot variant effect on deep mutational scans. Section~\ref{sec:tension}
argues that the two results are consistent and that the difference is diagnostic,
since a fine-tuned readout can recover residue identity from nucleotides through the
codon channel of Section~\ref{sec:map} while a likelihood cannot.

\paragraph{Evaluation critiques.} Our position follows a line of work showing that
benchmark gains often survive removal of the mechanism they are credited to
\citep{rendle2019difficulty,musgrave2020metric}. The specific concern that strong
scores reflect data leakage through a known map, and not learned structure, is
familiar from natural language inference artifacts \citep{gururangan2018annotation}.
We contribute the observation that biology supplies an exactly known map and
therefore exactly constructible controls, which natural language does not.

\paragraph{Noise ceilings.} Attenuation corrections are standard in psychometrics
and in neural encoding models, where repeat measurements bound explainable
variance. Our variant obtains the ceiling from model-internal invariance violation
and not from experimental replicates, which makes it available when replicates
are not.

\section{A Validity Protocol}

Let $f$ be a model scoring inputs in representation $A$, and let a claim assert
that $f$ captures property $P$ defined on modality $B$, related to $A$ by a known
surjective map $\pi: A \to B$. Four checks follow.

\subsection{A noise ceiling from model-internal invariance}
\label{sec:ceiling}

Because $\pi$ is many-to-one, a fibre $\pi^{-1}(b)$ is a set of inputs that must
share the label $P(b)$ exactly. Any variation of $f$ within a fibre is variance
that cannot track the label. Writing $\sigma^2_{\text{within}}$ for the mean
within-fibre variance of $f$ and $\sigma^2_{\text{between}}$ for the variance
across fibres, the correlation between $f$ and any label is attenuated by at most
\begin{equation}
\kappa \;=\; \sqrt{\frac{\sigma^2_{\text{between}}}{\sigma^2_{\text{between}} +
\sigma^2_{\text{within}}}}.
\end{equation}
A true correlation $\rho$ can appear no larger than $\kappa\rho$. Reporting a null
without $\kappa$ leaves open that the measurement had no power.

For genomic models the fibres are synonymous encodings. Crucially this costs
nothing: scoring a variant requires choosing a codon, and the same masked pass
assigns a probability to every synonymous alternative, so the within-fibre spread
is already computed. Across our 25 assays $\kappa = 0.900$, meaning a true
$\rho=0.5$ would still have surfaced at $0.45$.

\subsection{Map-destroying controls}
\label{sec:map}

A claim that $f$'s representation is $B$-level predicts that destroying $\pi$
destroys the signal. The control must destroy $\pi$ while keeping the input inside
$f$'s training distribution, or a drop is confounded with distribution shift.

For DNA two manipulations do this. A frameshift by one nucleotide changes every
downstream codon while leaving the string a real genomic sequence read one base
later. The reverse complement is real sequence on the opposite strand. Both destroy
the protein reading; neither makes the input unnatural.

We stress one negative design lesson. Replacing \emph{all} codons with synonymous
alternatives also preserves the protein, and its signal does drop, but this is a
global distribution shift as much as a map manipulation and is not admissible
evidence. The decisive controls are those that keep the input natural.

\subsection{Composition-matched context twins}
\label{sec:context}

Claims that additional context improves a score require a twin that supplies the
same number of tokens with the same composition and no information. Mononucleotide
and dinucleotide shuffles, applied separately to each flank so that each side's
composition is preserved, give this twin. The reported quantity is the real-minus-
shuffled difference, never the real trend alone.

\subsection{Trivial baselines and correlated units}
\label{sec:trivial}

Every table carries a scorer that uses no learned model. For protein fitness that
is BLOSUM62; for genomic benchmark tasks a positional one-hot or a $k$-mer count;
for representational alignment a one-hot residue encoding. Units are counted at the
level of independence: 194 ProteinGym assays constitute 165 protein clusters, and
our 25 assays constitute 21. Statistics are computed over clusters.

\section{Experimental Setup}

\paragraph{Assays.} Every ProteinGym assay for which a native coding sequence could
be verified, where verification requires that some ENA coding sequence for the
assay's UniProt accession translate to a protein containing the assayed target
exactly. After removing targets longer than 420 aa or with fewer than 200 usable
single mutants, 25 assays across 9 organisms remain: human, yeast,
\textit{E.\ coli}, \textit{B.\ subtilis}, \textit{P.\ aeruginosa},
\textit{K.\ pneumoniae}, \textit{S.\ acidocaldarius}, SARS-CoV-2 and HIV-1.

\paragraph{Scoring.} Identical masked-marginal protocol on both modalities, on the
same region: the assayed target sequence for protein models, and the coding
sequence encoding exactly it for genomic models. Only the modality differs. Evo~2
is autoregressive and is scored by summed token log-likelihood.

\paragraph{Models.} ESM-2 650M; Nucleotide Transformer v2 at 50M, 100M, 250M and
500M; HyenaDNA; Evo~2 7B. Genomic model code requires
\texttt{transformers==4.44.2}, since version 5 removes attributes the Nucleotide
Transformer's bundled remote code reads.

\section{Results}

\subsection{What the controls remove}

\paragraph{The transfer claim does not survive its map control.} A linear probe on
mean-pooled genomic embeddings, trained and evaluated within background groups so
that it must generalise to unseen residues, reaches $\rho = +0.141$. Read alone this
is weak transfer. Table~\ref{tab:map} is the same probe under manipulation.

\begin{table}[t]
\caption{A probe on genomic embeddings is invariant to destroying the reading
frame. 3{,}814 background groups, within-background $\rho$.}
\label{tab:map}
\centering\small
\begin{tabular}{llc}
\toprule
condition & protein & $\rho$ \\
\midrule
canonical & reference & $+0.141$ \\
random mutant codon & identical & $+0.143$ \\
frameshift $+1$ & \textbf{destroyed} & $+0.143$ \\
reverse complement & \textbf{destroyed} & $+0.129$ \\
\bottomrule
\end{tabular}
\end{table}

Frameshifting changes every downstream codon and costs nothing. Reading the
opposite strand, where the sequence is not a coding sequence at all, costs $0.012$.
No protein-level representation can be invariant to either. The mechanism is that
the codon at the varying site already determines the residue, so the probe never
needed a protein-level representation.

A variance decomposition agrees. The fraction of score variance explained by
amino-acid identity is $0.643$ for the genomic likelihood and $0.672$ for the
genomic representation, against $1.000$ for any protein model by construction. The
representation is barely more protein-level than the likelihood.

\paragraph{A genomic model aligns with ESM-2 worse than a lookup table does.}
Linear CKA and a cross-modal ridge map with whole amino acids held out:

\begin{table}[h]
\centering\small
\begin{tabular}{lcc}
\toprule
pair & CKA & held-out $R^2$ \\
\midrule
NT $\rightarrow$ ESM-2 & $0.663$ & $+0.412$ \\
\textbf{one-hot} $\rightarrow$ ESM-2 & $\mathbf{0.819}$ & $\mathbf{+0.553}$ \\
\bottomrule
\end{tabular}
\end{table}

A 60-parameter one-hot vector shares more with ESM-2 than the genomic model does.

\paragraph{Context trends reproduce from no information.} For the Nucleotide
Transformer, real genomic flanks produce a monotone degradation with trend
$-0.771$; mononucleotide-shuffled flanks reproduce it at $-0.633$ over three seeds.
HyenaDNA with a 30~kb window produces a clean positive trend of $+0.536$ whose
shuffled twins score $+0.893$, $+0.750$ and $+0.750$, so real minus shuffled is
$-0.008$ and real is higher in 1 of 6 comparisons. Both signs are reproduced by
flanks that carry nothing.

For Evo~2 the control is different and sharper. No arm of its context curve carries
specificity signal at all: the bare coding sequence scores AUC $0.401$ with
interval $[0.345,0.456]$, significantly \emph{below} chance, and every flanked arm
lies between $0.521$ and $0.543$ with an interval containing $0.5$. The reported
$+0.70$ trend over five points carries an exact permutation $p$ of $0.117$, falls
to $+0.400$ with $p=0.375$ once the no-context point is dropped, and the 300~nt and
5400~nt score vectors correlate at $+0.963$. An eighteen-fold increase in context
past the first increment changes AUC by $+0.016$ and then nothing.

\subsection{Transfer across 25 assays}

\begin{table}[t]
\caption{Mean Spearman $\rho$ across 25 DMS assays, identical protocol and region.
Intervals are 95\% over assays.}
\label{tab:main}
\centering\small
\begin{tabular}{lcc}
\toprule
scorer & mean $\rho$ & 95\% CI \\
\midrule
ESM-2 650M (protein) & $+0.4655$ & $[+0.395,+0.536]$ \\
chemistry ridge, LOCO & $+0.2860$ & $[+0.248,+0.324]$ \\
\textbf{Evo~2 7B (genomic)} & $\mathbf{+0.2658}$ & $[+0.202,+0.330]$ \\
BLOSUM62 alone & $+0.2282$ & $[+0.195,+0.261]$ \\
NT-v2 50M (genomic) & $-0.0132$ & $[-0.033,+0.007]$ \\
\bottomrule
\end{tabular}
\end{table}

Table~\ref{tab:main} is the headline. Three readings follow, and the trivial
baseline decides all three.

A substitution matrix with no model of any kind reaches 49\% of ESM-2 and is
positive on 24 of 25 assays. Five fitted chemistry scalars with weights learned on
other protein clusters reach 61\%. ESM-2's genuinely learned margin over fitted
chemistry is $+0.179$, not $+0.465$. The bar a genomic model must clear is
therefore $+0.228$, not $+0.466$.

The Nucleotide Transformer does not clear it, losing to BLOSUM62 by $+0.241$
$[+0.200,+0.283]$ in 24 of 25 assays. Its attenuation ceiling is $0.900$, so this
is a null and not a power failure. Marginalising over synonymous codons, which
removes the leak of Section~\ref{sec:map} by construction and is the only genomic
score entitled to be called a protein-level proxy, scores $-0.019$.

Evo~2 does clear zero, at $+0.266$ with an interval excluding it and positive on 23
of 25 assays. It does not clear the substitution matrix: Evo~2 minus BLOSUM62 is
$+0.0376$ with interval $[-0.012,+0.087]$, higher in 14 of 25. Against the fitted
chemistry ridge it is behind, at $-0.0202$.

\subsection{Scale is not the explanation}
\label{sec:scale}

The Nucleotide Transformer is flat across four sizes. Fitting $\rho$ against
$\log_{10} N$ on the four points gives a slope of $+0.0179$ per decade and an
intercept of $-0.1516$, with Spearman $+0.800$ over four points. Extrapolated, that
trend predicts $+0.0241$ at 7B parameters against Evo~2's measured $+0.2658$, short
by a factor of $11.0$. Reaching BLOSUM62 on that trend would require
$10^{21.3}$ parameters.

Evo~2 therefore sits far off the line that scaling its family would trace. Whatever
produces its protein-fitness signal is a property of architecture and corpus, and
scale alone does not buy it. We note the limitation squarely: with one point off
the line we can reject the scaling account but cannot identify which of
architecture, corpus or token budget is responsible, because none of these is a
variable we control. Section~\ref{sec:future} says what would.

\subsection{Architecture and corpus, measured from released checkpoints}
\label{sec:stage0}

Section~\ref{sec:scale} rejects the scaling account but cannot say which of
architecture, corpus or token budget produces Evo~2's signal, because those are
confounded between the two families. Released checkpoints resolve two of the three
at matched architecture and inference cost only. We ran both.

\paragraph{Codon-aligned tokenization does not help.} The v2 family ships a 3-mer
variant that is otherwise identical to the 6-mer 50M model: 512 hidden, 12 layers,
16 heads, the same 850-genome corpus, vocabulary 4107 against 75. Because a 3-mer
token is exactly a codon, the masked token \emph{is} the codon under test and the
neighbouring-codon conditioning the 6-mer scorer carries disappears.
\citet{boshar2024glm} trained this model because codon-level resolution was where
genomic models lost to protein models under fine-tuning, and reported that it
improved their two codon-sensitive tasks. It had not been scored zero-shot.

\paragraph{Within-species corpus diversity does not help.} The v1 family supplies
\texttt{500m-human-ref} and \texttt{500m-1000g}, identical at 1280 hidden, 24 layers
and 20 heads, differing only in training on one human reference genome against 3,202
genetically diverse human genomes.

\begin{table}[t]
\caption{Contrasts available from released checkpoints, each varying one factor at
matched architecture. Paired bootstrap over 25 assays. Every difference contains
zero, and every model misses the BLOSUM62 bar by about $0.23$.}
\label{tab:free}
\centering\small
\begin{tabular}{lcccc}
\toprule
model & mean $\rho$ & 95\% CI & $\kappa$ & vs.\ B62 \\
\midrule
\multicolumn{5}{l}{\emph{tokenization, 850-genome corpus, 50M}}\\
\quad 6-mer & $-0.0132$ & $[-0.033,+0.006]$ & $0.900$ & $-0.241$ \\
\quad 3-mer & $-0.0288$ & $[-0.053,-0.003]$ & $0.866$ & $-0.257$ \\
\multicolumn{5}{l}{\quad paired: $+0.0156\ [-0.012,+0.043]$, contains zero}\\
\addlinespace
\multicolumn{5}{l}{\emph{corpus diversity, 500M}}\\
\quad human-ref & $+0.0032$ & $[-0.015,+0.021]$ & $0.901$ & $-0.225$ \\
\quad 1000G & $-0.0022$ & $[-0.024,+0.017]$ & $0.878$ & $-0.230$ \\
\multicolumn{5}{l}{\quad paired: $+0.0054\ [-0.013,+0.025]$, contains zero}\\
\bottomrule
\end{tabular}
\end{table}

Table~\ref{tab:free} reports both. Neither factor moves the result. Codon-aligned
tokenization gives a paired difference of $+0.0156$ $[-0.012,+0.043]$, and three
thousand genomes of within-species variation give $+0.0054$ $[-0.013,+0.025]$. Both
contain zero. All four models sit about $0.23$ below the substitution-matrix bar,
and every attenuation ceiling lies between $0.866$ and $0.901$, so none of these
nulls is a power failure.

We note one sampling lesson from the run itself. At 14 of the 25 assays the
tokenizer difference read $-0.0302$ with an interval excluding zero, and that did
not survive the remaining eleven. These intervals are wide at small $n$ and a
partial sweep is not a result.

\paragraph{What remains unmeasured.} Cross-species corpus breadth, available as
NT~2.5B trained on 1000G against the same size trained on 850 species, is the one
free contrast we could not run, for storage and not compute reasons. It is also the
most informative of the three, since Evo~2 trains across all domains of life and
breadth is the most plausible remaining account of its margin. We report its absence
explicitly because the conclusion we draw in Section~\ref{sec:future} would change
if it showed a gradient.

\subsection{Genomic models on their own task}

Grading genomic models on coding sequence and protein fitness is grading them on a
task they were not built for. On five Nucleotide Transformer benchmark tasks under
a homology-clustered split, with frozen embeddings and a linear probe, the trivial
competitor wins outright twice and loses twice by under $0.03$ AUC. A positional
one-hot reaches $0.959$ on splice sites against the Nucleotide Transformer's
$0.756$ and HyenaDNA's $0.704$, a loss of $0.203$ and $0.254$. The one-hot
competitor is selected on test, so the baseline is given an oracle the model is
not, and still wins.

\section{Reconciling Fine-Tuned and Zero-Shot Results}
\label{sec:tension}

\citet{boshar2024glm} report Nucleotide Transformer v2 500M matching or beating
ESM-2 on three of five protein tasks. We report the same family at $\rho=-0.013$
against ESM-2's $+0.466$. Both can hold, and the protocol says why.

Fine-tuning fits a readout on top of the representation. Section~\ref{sec:map}
shows that the genomic representation carries residue identity through the codon,
and that this survives frameshifting and reverse-complementing. A supervised head
is free to use that channel, so strong fine-tuned performance on a protein task is
evidence that residue identity is recoverable and not that the model represents
protein-level structure. A zero-shot likelihood has no fitted head and therefore no
access to the channel, which is why the same model scores at chance.

This makes the two measurements complementary and gives a testable prediction: a
model whose fine-tuned protein performance comes from the codon channel should show
that performance survive a frameshift of its input. We did not run this and regard
it as the sharpest available test of the account, since it requires only re-running
an existing fine-tuning pipeline on frameshifted sequences.

A codon-aligned tokenization is the natural place to look for a genuine effect, and
one exists at matched architecture: \texttt{nucleotide-transformer-v2-50m-3mer-}
\texttt{multi-species} differs from its 6-mer counterpart only in vocabulary size,
75 against 4107, with identical depth, width, head count and training corpus. We
report it zero-shot in Table~\ref{tab:free} as the one architectural change prior
work identified as promising for protein tasks.

\section{Four Claims the Protocol Overturned}
\label{sec:retract}

\textbf{A reading-frame result was retracted.} A genomic model appeared unable to
distinguish a gene from the same gene read one nucleotide out of frame. The
``frameshift'' is a rotation, so apart from one wrap junction it is real genomic
sequence. Ranking the four conditions by whether they yield a real genomic string
reproduces the entire table with no appeal to reading frame. The replacement probe,
calibrated against a 3-periodic Markov positive control, is underpowered and not
null: bounding it needs about 214 unique coding sequences and 24 are available.

\textbf{A representation-level transfer claim was falsified by its own control},
as Table~\ref{tab:map} shows.

\textbf{A context dose-response was measuring composition}, in both directions.

\textbf{A general claim was too general.} ``A genomic model carries no
protein-fitness signal'' held for the Nucleotide Transformer family and was written
as though it held for genomic models. Evo~2 refutes it.

\section{What the Protocol Costs}

Adoption is cheap, which is the argument for it.

The noise ceiling is free: it is read off a masked pass already performed. The map
controls cost one additional scoring pass per manipulation, so two passes over the
evaluation set. Composition-matched twins cost one pass per shuffle seed per
condition; we use two seeds and two shuffle orders. Trivial baselines cost
negligible compute. Cluster-level statistics cost nothing but a join.

For the Evo~2 context control specifically, scoring the 399-variant discrimination
subset instead of all 7{,}882 variants is a twentyfold saving that changes no
reported number, because the published AUCs are read off that subset already. The
full control is 21 arms, 8{,}379 sequences and 41.9M tokens, against 167M for the
original uncontrolled arms. A control that is cheaper than the experiment it checks
has no compute excuse.

\section{Limitations}
\label{sec:future}

Evo~2 numbers come from a single run on one GPU host and describe \texttt{evo2\_7b},
since \texttt{evo2\_1b\_base} requires Transformer Engine. We have one genomic
architecture off the Nucleotide Transformer scaling line and therefore cannot
separate architecture from corpus from token budget.

Public releases fix the \emph{token budget} at one value per model: we verified
that neither the Nucleotide Transformer nor Evo~2 publishes intermediate-step
checkpoints, every repository carrying a single branch and one weight file. The
token axis is therefore closed to any study of frozen models.

Corpus and tokenization are not closed, and Section~\ref{sec:stage0} measures
both. Neither moves the result, which bears on what a pretraining ladder could be
expected to find. Of the three factors separable without training a model, two are
flat and the third is unmeasured here.

What remains genuinely unobservable from released models is the token budget and
any factor crossed with it. A ladder pretrained under controlled variation of
parameters, tokens and corpus composition would close it, and Appendix~A gives a
design and budget. We state the case against that ladder as plainly as the case
for it: with two of three free axes null, the argument for the spend rests on the
token axis alone, and we would not undertake it before the breadth contrast of
Section~\ref{sec:stage0} is run, since that is three orders of magnitude cheaper
and bears on the same question.

Every assay used here has homologs, so the regime where a sequence has no relatives
is untested. The regulatory arm covers two small models and licenses nothing about
larger ones. The specificity arm rests on one combinatorially complete landscape.

\section{Conclusion}

The controls in this paper are not specific to biology. Any claim that a
representation encodes a property of another modality related to it by a known map
can be checked by destroying the map, by bounding attenuation from within-fibre
variance, and by matching a context manipulation on composition. Biology is the
setting where the map is known exactly, which makes it the setting where the
controls are exactly constructible and where their consequences can be seen most
clearly. In this setting they removed four claims, three of them ours, and left a
smaller result standing: genomic-to-protein transfer exists at Evo~2 scale and does
not clear a substitution matrix. Two factors that might have explained it, codon
tokenization and corpus diversity, do not, and we were able to establish that from
released checkpoints without training anything. A reviewer entitled to ask what
would change our mind has a short answer: cross-species corpus breadth, at a cost
of one download.

\bibliographystyle{icml2027}
\bibliography{refs}

\end{document}
